\section{GALA Code Construction}\label{sec:construction}
Let's consider another generalization of the block circulant codes where the non-commutativity is instead supplied by a small non-abelian finite group.


\subsection{Monomial Direct Product Codes}\label{ssec:mono-gala}
\begin{definition}[Direct Product GALA Code]\label{def:dpg-code}
Given parameters $L, J\leq \tfrac L2, k, m\in\mathbb Z_+$ and an active set $\Gamma$ as in Definition~\ref{def:kasai}, the Direct Product GALA (DPG) codes $\mathrm{GALA}_{L,J}(H_k\times C_m)$ are the family of two-block CSS codes on $n = Lkm$ qubits lifted over the group algebra $\mathbb F_2[H_k\times C_m]$, where $H_k$ is a non-abelian permutation group on $[k]$ and $C_m$ is an abelian permutation group on $[m]$. A DPG code is specified by choosing the generators (also called lifts) $\mathcal{F},\mathcal{G}\subset H_k\times C_m$:
\begin{align*}
    \mathcal{F} &= \{(F_i, f_{i})\}_{i\in[\frac L2]},\\
    \mathcal{G} &= \{(G_i, {g}_i)\}_{i\in[\frac L2]},
\end{align*}
with $F_i, G_i\in H_k$ and $f_i,g_i\in C_m$, such that
\begin{equation}\label{eq:dpg-active}
    [F_i, G_j] = 0\iff (i,j)\in \Gamma .
\end{equation}
By Lemma~\ref{lemma:dpk-commutivity}, $[(F_i,f_i),(G_j,g_j)] = 0\iff [F_i,G_j] = 0$, so \eqref{eq:dpg-active} is exactly the commutativity pattern of Definition~\ref{def:kasai} for the lifted generators: the entire orthogonality pattern is decided inside $H_k$, and the abelian entries $f_i,g_i$ are unconstrained by it. The stabilizers $ H_X, {H}_Z$ are given by the first $J$ block rows of the parent matrices $\hat H_X = [F|G]$ and $\hat H_Z = [G^T|F^T]$, where
\begin{align*}
    F &= \sum_{i\in [\frac{L}{2}]} z_i\otimes F_i\otimes f_i ,\\
    G &= \sum_{i\in [\frac{L}{2}]} z_i\otimes G_i\otimes g_i,
\end{align*}
so that $F,G\in \mathbb F_2[\mathbb Z_{\frac L2}\times H_k\times C_m]$, with $z_i\in\mathbb Z_{\frac L2}$ as in Definition~\ref{def:kasai}.
\end{definition}

Notice that when using a small group $H_k$ as the non-abelian part (say $S_3$), it is computationally simple to enumerate all tuples $F_i, G_i\in H_k$ satisfying a given active orthogonality constraint \eqref{eq:dpg-active}, and the enumeration is independent of $m$. Two remarks make this enumeration a finite classification rather than a search. First, only the tops enter \eqref{eq:dpg-active}. Second, the classification need only be carried out up to the symmetries of the construction, of which the following reflection is the one we use.

Let us give an explicit example. Suppose $L = 12$, $J = 3$, pick the active set $\Gamma = \left([\frac L2]\times [\frac L2]\right) \setminus \{(0,3), (1,2)\}$ of Definition~\ref{def:kasai}, and let $H_k = S_3$. We can write down this group and its multiplication table as follows:

$$
\renewcommand{\arraystretch}{1.3}
\begin{array}{c|cccccc}
\cdot      & e        & \sigma_0 & \sigma_1 & \tau_0   & \tau_1   & \tau_2   \\ \hline
e          & e        & \sigma_0 & \sigma_1 & \tau_0   & \tau_1   & \tau_2   \\
\sigma_0   & \sigma_0 & \sigma_1 & e        & \tau_2   & \tau_0   & \tau_1   \\
\sigma_1   & \sigma_1 & e        & \sigma_0 & \tau_1   & \tau_2   & \tau_0   \\
\tau_0     & \tau_0   & \tau_1   & \tau_2   & e        & \sigma_0 & \sigma_1 \\
\tau_1     & \tau_1   & \tau_2   & \tau_0   & \sigma_1 & e        & \sigma_0 \\
\tau_2     & \tau_2   & \tau_0   & \tau_1   & \sigma_0 & \sigma_1 & e        \\
\end{array},
\qquad
\begin{array}{c|c}
\text{element} & \text{cycle} \\ \hline
e        & \mathrm{id}   \\
\sigma_0 & (0\,1\,2)     \\
\sigma_1 & (0\,2\,1)     \\
\tau_0   & (1\,2)        \\
\tau_1   & (0\,2)        \\
\tau_2   & (0\,1)        \\
\end{array}.
$$
Commutativity can be checked by a table look-up, and the possible tops are then completely determined.

\begin{lemma}[$S_3$ ansatzes at $L=12$, $J=3$]\label{lemma:s3-ansatz}
Let $L=12$, $J=3$, $\Gamma = \left([\frac L2]\times [\frac L2]\right)\setminus\{(0,3),(1,2)\}$ and $H_k = S_3$. Then the tuples $(\mathcal F^H,\mathcal G^H)$ satisfying \eqref{eq:dpg-active} are exactly
\begin{align}
    \mathcal{F}^H &= \{\sigma_j,\ \tau_i,\ e,\ e,\ e,\ e\},
        & \mathcal{G}^H &= \{e,\ e,\ \sigma_k,\ \tau_i,\ e,\ e\};\label{eq:ansatz-a}\\
    \mathcal{F}^H &= \{\tau_i,\ \sigma_j,\ e,\ e,\ e,\ e\},
        & \mathcal{G}^H &= \{e,\ e,\ \tau_i,\ \sigma_k,\ e,\ e\};\label{eq:ansatz-b}\\
    \mathcal{F}^H &= \{\tau_i,\ \tau_j,\ e,\ e,\ e,\ e\},
        & \mathcal{G}^H &= \{e,\ e,\ \tau_i,\ \tau_j,\ e,\ e\}, \quad i\neq j,\label{eq:ansatz-c}
\end{align}
\end{lemma}

\begin{proof}
It is easy to check exhaustively.
\end{proof}

To get a valid code we then choose the bottom entries $f,g\in C_m$ greedily to optimize girth and distance, which by Lemma~\ref{lemma:dpk-commutivity} costs nothing in orthogonality. The $[[1152,580,\leq12]]$ code of \cite{zhao2026ultrahighratequantumerrorcorrection} is an instance of this example with $k = 3$, $H_3 = S_3$ and $C_m = \mathbb Z_{32}$, so that $n = Lkm = 12\cdot3\cdot32 = 1152$. 


Before proving this connection, and the connections to the other Kasai codes, let us record the more general version of the GALA codes:
\begin{definition}[Semidirect Product GALA Code]\label{def:spg-code}
Given parameters $L, J \leq \tfrac{L}{2}, k, m \in \mathbb{Z}_+$ and an active set $\Gamma$ as in Definition~\ref{def:kasai}, the Semi-direct Product GALA (SPG) codes are the family of two-block CSS codes on $n = Lkm$ qubits lifted over the group algebra $\mathbb F_2[H_k\ltimes C^k_m]$, where $H_k$ is a non-abelian permutation group on $[k]$, $C_m$ is an abelian permutation group on $[m]$, and $H_k\ltimes C_m^k$ is the semi-direct product of Definition~\ref{def:semidirect-product}. An SPG code is specified by choosing the generators $\mathcal{F},\mathcal{G}\subset H_k\ltimes C^k_m$:
\begin{align*}
    \mathcal{F} &=  \{(F_i;\textbf{f}_i )\}_{i\in[\frac L2]},\\
    \mathcal{G} &= \{(G_i; \textbf{g}_i)\}_{i\in[\frac L2]},
\end{align*}
such that $[(F_i;\textbf{f}_i),(G_j;\textbf{g}_j)] = 0\iff (i,j)\in\Gamma$, which by Lemma~\ref{lemma:spk-commutivity} is the commutativity condition~\eqref{eq:spk-commutivity} on the pair $(F_i;\textbf f_i), (G_j;\textbf g_j)$. The stabilizers $ H_X, {H}_Z$ are given by the first $J$ block rows of the parent matrices $\hat H_X = [F|G]$ and $\hat H_Z = [G^T|F^T]$, where
\begin{align*}
    F &= \sum_{i\in [\frac{L}{2}]} z_i\otimes M_\ltimes(F_{i}, \textbf{f}_{i}) ,\\
    G &= \sum_{i\in [\frac{L}{2}]} z_i\otimes M_\ltimes(G_{i},\textbf{g}_{i}).
\end{align*}
\end{definition}

The commutativity condition~\eqref{eq:spk-commutivity} is harder to work with than \eqref{eq:dpg-active}, since the bottoms no longer drop out; but Lemma~\ref{lemma:spk-commutivity} still supplies simplifying sufficient conditions, which serve as good heuristics for computing ansatzes and simplify the proofs below. The most useful of these is to make the bottoms constant on a subset of the lifts:

\begin{definition}[Restricted SPG Code]\label{def:SPG-restricted}
A restricted semidirect Product GALA code $\mathrm {GALA}_{L,J}(H_k\ltimes_\mathcal{R} C_m)$ on $\mathcal{R}\subseteq [L]$ is an SPG code with  $\textbf{f}_i = \textbf{1}\otimes f_i$, $f_i\in C_m$ for all $i\in \mathcal{R}$ and $\textbf{g}_j = \textbf{1}\otimes g_j$, $g_j\in C_m$  for all  $j + \frac L2\in \mathcal{R}$.
\end{definition}

If we restrict an SPG code on the maximal $\mathcal{R} = [L]$, all its matrix elements simplify to $M_\ltimes(h, \textbf{1}\otimes c) = M_\times(h\otimes c)$; that is, the DPG codes are exactly the SPG codes with maximal restriction. The size of the restriction corresponds exactly to the number of abelian column generators in \cite{zhao2026ultrahighratequantumerrorcorrection}, and it is what controls the cost of the movement schedule in Section~\ref{sec:hardware}.

\subsection{Permutation Equivalence of Kasai}\label{ssec:kasai-perm-eq}
We can now show that GALA codes contain Kasai codes as special cases, once AOD compatibility is required of the latter. We say two codes $\mathcal{P}, \mathcal{Q}$ with check matrices $H_P, H_Q$ are permutation equivalent, or $\mathcal{P}\sim_\pi \mathcal{Q}$, when $\pi_r\cdot H_P\cdot \pi_c = H_Q$ for some permutation matrices $\pi_r, \pi_c$. For quantum codes this condition must hold for both $H_X$ and $H_Z$ with the same $\pi_c$. It guarantees that the two codes have the same logical space up to a relabeling of the qubits, and in particular the same $n,k,d$ and girth.

\begin{proposition}[Kasai codes with a reference APM are GALA codes]\label{prop:kasai-containment}
Let $\mathcal{K}=\mathrm K_{L,J}(\mathcal F,\mathcal G)$ be a Kasai code on $[P]$ with generators $\mathcal{F},\mathcal{G}\subset\mathbb A_P$ and active set $\Gamma$ as in Definition~\ref{def:kasai}.
\begin{enumerate}
    \item \emph{(Semi-direct.)} Suppose the centralizer $\mathcal A = Z_{\mathbb A_P}(\mathcal F\cup\mathcal G)$ contains an element $A$ acting \emph{freely} on $[P]$ such that $A^s$ has no fixed point for $0<s<\ord{A}$, and put $m = \ord A$. Then $k = P/m$ is an integer and $\mathcal K\sim_\pi\mathcal Q$ for some $\mathcal Q\in\mathrm{GALA}_{L,J}(H_k\ltimes\mathbb Z_m)$ with the same $L,J,\Gamma$, where $\mathbb Z_m = \rs{A}$ acts regularly inside each orbit of $A$ and $H_k\leq S_k$ is the group induced by $\mathcal F\cup\mathcal G$ on the set of $k$ orbits.
    \item \emph{(Direct.)} Suppose instead $P = km$ with $\gcd(k,m) = 1$, and that the reductions modulo $m$ of the generators pairwise commute in $\mathbb A_m$. Then $\mathcal K\sim_\pi\mathcal Q$ for some $\mathcal Q\in\mathrm{GALA}_{L,J}(H_k\times C_m)$ with the same $L,J,\Gamma$, where $H_k\leq\mathbb A_k$ and $C_m\leq\mathbb A_m$ are generated by the reductions of $\mathcal F\cup\mathcal G$ modulo $k$ and modulo $m$ respectively.
\end{enumerate}
\end{proposition}

\begin{proof}
The first case follow from Lemma 1 of~\cite{zhao2026ultrahighratequantumerrorcorrection}, which we will restate in our language. Let us fix orbit representatives $\gamma = \gamma_0,\dots,\gamma_{k-1}$ and relabel $[P]$ by
\[
    \pi : [k]\times[m]\to[P],\qquad (a,s)\mapsto A^s\gamma_a ,
\]
which is bijective since $A$ act freely on $P$. Let $M\in Z_{\mathbb A_P}(A)$. Then $M\gamma_a = A^{s_a}\gamma_{h(a)}$ for some $h(a)\in[k]$ and $s_a\in\mathbb Z_m$, and $MA^s = A^sM$ gives
\[
    M(A^s\gamma_a) = A^sM\gamma_a = A^{s+s_a}\gamma_{h(a)};
\]
that is, $M$ maps each orbit of $A$ onto an orbit of $A$ and acts inside it as a uniform shift, which is exactly the wreath product structure where each orbit correspond to a shift on $[m]\cong \gamma$ and the inter-orbit shifts correspond to a permutation on $[k]$. The second case follows from CRT.

\end{proof}


Indeed, we can verify numerically that the $[[1152,580]]$ code is permutation equivelent to $\mathrm{GALA}_{12,3}(S_3\times \mathbb Z_{32}) $, and the $[[2304,1156]]$ codes are permutation equivalent to $\mathrm{GALA}_{12,3}(S_3\times\mathbb Z_2\times  \mathbb Z_{32})$.

As we show in Section~\ref{sec:bounds}, the monomial DPG construction that covers the co-designed Kasai codes~\cite{zhao2026ultrahighratequantumerrorcorrection} faces fundamental limits on the achievable parameters: reaching $d\geq L$ at small $n$, and reaching higher $d$ at all, requires further generalizations. We give these generalizations here and discuss how they overcome the barriers in Section~\ref{sec:bounds}.

\subsection{Polynomial and Loose-Active Orthogonality Generalizations}\label{ssec:generalized-gala}

Both constructions so far lift each proto-matrix entry to a single group element. Upgrading some lifts to a \emph{sum} of group elements --- that is, taking the lift in the group ring $\mathbb F_2[H_k\star C_m]$ rather than in the group $H_k\star C_m$ --- enlarges the search space at fixed $L$, and in particular by-passes the distance bound~\eqref{eq:j<=2} of Proposition~\ref{prop:distance-bound}, so that $L = 8$, $J = 2$ codes become available. Once the lifts are sums, a second relaxation follows for free: the term-wise condition $[F_{i,i'},G_{j,j'}] = 0$ for $(i,j)\in\Gamma$ is stronger than active orthogonality requires, since by~\eqref{eq:psi-r} what must vanish is the aggregate $\Psi_r$, and the individual commutators of a polynomial lift can cancel against each other. This gives the final generalization of loose-active orthogonality.


\begin{definition}[Polynomial GALA Code]\label{def:pg-code}
Given parameters $L, J \leq \tfrac{L}{2}, k, m \in \mathbb{Z}_+$ and an active set $\Gamma$ as in Definition~\ref{def:kasai}, let $G = H_k\star C_m$ be either the direct product $H_k\times C_m$ of Definition~\ref{def:direct-product} or the semi-direct product $H_k\ltimes C^k_m$ of Definition~\ref{def:semidirect-product}, with $H_k$ a non-abelian permutation group on $[k]$ and $C_m$ an abelian permutation group on $[m]$. The Polynomial GALA codes $\mathrm{GALA}_{L,J}(\mathbb F_2[G])$ are the family of two-block CSS codes on $n = Lkm$ qubits in which each lift entry is an element of the \emph{group ring} $\mathbb F_2[G]$, i.e.\ a sum of group elements rather than a single one. A polynomial GALA code is specified by choosing generators $\mathcal F,\mathcal G\subset\mathbb F_2[G]$,
\begin{align*}
    \mathcal{F} &= \left\{\sum_{i'\in I^{\mathcal F}_i}(F_{i,i'}; \textbf{f}_{i,i'})\right\}_{i\in[\frac L2]},\\
    \mathcal{G} &= \left\{\sum_{j'\in I^{\mathcal G}_j}(G_{j,j'}; \textbf{g}_{j,j'})\right\}_{j\in[\frac L2]},
\end{align*}
over finite index sets $I^{\mathcal F}_i, I^{\mathcal G}_j$, such that every pair of monomials $(F_{i,i'};\textbf f_{i,i'})$, $(G_{j,j'};\textbf g_{j,j'})$ with $(i,j)\in\Gamma$ commutes --- by Lemma~\ref{lemma:dpk-commutivity} in the direct-product case and by condition~\eqref{eq:spk-commutivity} of Lemma~\ref{lemma:spk-commutivity} in the semi-direct case --- while some pair with $(i,j)\notin\Gamma$ does not. The stabilizers $ H_X, {H}_Z$ are given by the first $J$ block rows of the parent matrices $\hat H_X = [F|G]$ and $\hat H_Z = [G^T|F^T]$, where
\begin{align*}
    F &= \sum_{i\in [\frac{L}{2}]} z_i\otimes \sum_{i'\in I^{\mathcal F}_i}M_\ltimes(F_{i,i'}, \textbf{f}_{i,i'}) ,\\
    G &= \sum_{i\in [\frac{L}{2}]} z_i\otimes \sum_{i'\in I^{\mathcal G}_i}M_\ltimes(G_{i,i'},\textbf{g}_{i,i'}),
\end{align*}
reading $M_\ltimes(h,\textbf 1\otimes c) = h\otimes c$ in the direct-product case, so that $F,G\in\mathbb F_2[\mathbb Z_{\frac L2}\times G]$. The monomial codes of Definitions~\ref{def:dpg-code} and~\ref{def:spg-code} are the case $|I^{\mathcal F}_i| = |I^{\mathcal G}_j| = 1$ for all $i,j$.
\end{definition}

\begin{definition}[Polynomial GALA Code with Loose-Active Orthogonality]\label{def:loose-orthogonality}
A polynomial GALA code with loose-active orthogonality is a code built as in Definition~\ref{def:pg-code} in which the term-wise commutation requirement on $\Gamma$ is replaced by the weaker aggregate requirement
\begin{equation}\label{eq:loose-active}
    \Psi_r = 0\quad\text{for every } r\in\{j-i\ :\ i,j\in[J]\},\qquad \Psi_{r'}\neq 0\ \text{ for some } r'\in[\tfrac L2],
\end{equation}
with $\Psi_r$ as in~\eqref{eq:psi-r}. Grouping the terms of $\Psi_r$ under the involution $u\mapsto r-u$, \eqref{eq:loose-active} is implied by the paired conditions
\[
    [\mathcal F_u,\mathcal G_{r-u}] + [\mathcal F_{r-u},\mathcal G_u] = 0\quad\text{for } u\in[\tfrac L2],
    \qquad\text{and}\qquad
    [\mathcal F_u,\mathcal G_u] = 0 \quad\text{whenever } 2u\equiv_{\frac L2} r,
\]
imposed for each active offset $r$, which may hold even when no individual commutator vanishes.
\end{definition}


\subsection{Codes with ZX-dualities}\label{ssec:zxdual}
Finally, we can impose even more symmetry during the code search, so as to improve parallelization and the availability of logical operations. The relevant symmetry here is a ZX-duality in the sense of Definition~\ref{def:zx-duality}, and for a two-block code it is available whenever the two lists of lifts are related pointwise by a re-indexing of the blocks.

\begin{proposition}[ZX-Dual GALA Codes]\label{def:dual-gala}
Let $\mathcal{Q}\in \mathrm{GALA}_{L,J}(H_k\star C_m)$ with $4\mid L$, and let $R = \{r_0,r_1, r_2, r_3 \}$ denote the following involutions on $[\frac{L}{2}]$:
\begin{equation}\label{eq:zx-sector-maps}
    r_0 = e:\; i\mapsto i,\qquad r_1:\; i\mapsto i + \tfrac L4,\qquad r_2:\; i\mapsto -i,\qquad r_3=r_1r_2:\; i\mapsto \tfrac L4 - i ,
\end{equation}
If the generators of $\mathcal Q$ satisfy
\begin{equation}\label{eq:zx-generator-condition}
    F_iG_{r(i)} = t \quad \forall i\in [\tfrac L2]
\end{equation}
for some {sector involution} $r\in R$, we say $\mathcal Q$ is ZX-dual. The ZX-duality is given by $\tau = \iota_r\otimes t$, where
\begin{equation}\label{eq:zx-iota}
\iota_{r_0} = r_3, \iota_{r_1} = r_2, \iota_{r_2} = r_0, \iota_{r_3} = r_1.
\end{equation}
\end{proposition}

\begin{proof}
The first two are \emph{translations} of $\mathbb{Z}_\frac L2$ and the last two are \emph{reflections}. It is easy to see that when $r =r_2$, $F_i = G_{-i}^T\cdot t$, and $H_X\cdot e\otimes t = H_Z$ exactly. It remains three cases. For $r = r_0$, $\iota = r_1r_2$. We can see this by tracing the matrix elements:
\begin{align*}
(H_X)_{i,j} 
&\xrightarrow{r_2} (H_X)_{i,-j}\\
&\xrightarrow{r_1} (H_X)_{i, -j+\frac L4} \\
&= F_{i-(-j+\frac L4)} \\
&=G^T_{i-(-j+\frac L4)}\cdot t \\
&=G^T_{j-(\frac L4-i)}\cdot t =(H_Z)_{\frac L4-i,j}. 
\end{align*}
The derivation for other cases are similar.
\end{proof}



A natural question one might ask is whether this restriction leads to girth or distance failures. It does, for the reflection sector involutions, and we return to it in Section~\ref{ssec:zx-bounds}.
